\chapter{Modelado UML y Diagramas ER}

\section{Diagrama Entidad-Relación (ER)}

Modelo conceptual independiente del DBMS. Define entidades, atributos y relaciones antes de pensar en tablas.

\begin{itemize}
  \item \textbf{Entidad:} objeto del mundo real con existencia propia (rectángulo)
  \item \textbf{Atributo:} propiedad de la entidad (elipse / columna en Crow's Foot)
  \item \textbf{Relación:} asociación entre entidades (rombo en Chen, línea en Crow's Foot)
  \item \textbf{Entidad débil:} depende de otra para existir (sin PK propia)
\end{itemize}

\section{Notaciones}

\begin{itemize}
  \item \textbf{Chen:} rectángulos, rombos, elipses. Académica, verbosa.
  \item \textbf{Crow's Foot:} líneas con símbolos de cardinalidad en los extremos. Estándar en herramientas (Lucidchart, dbdiagram.io).
  \item \textbf{UML Clases:} atributos y métodos en compartimentos. Útil cuando el modelo de datos y el de objetos son lo mismo.
\end{itemize}

\section{Cardinalidad en Crow's Foot}

\begin{itemize}
  \item \texttt{||} exactamente uno (obligatorio)
  \item \texttt{|o} cero o uno (opcional)
  \item \texttt{||<} uno o muchos
  \item \texttt{|o<} cero o muchos
\end{itemize}

\begin{diagramabox}{Diagrama ER --- Cliente, Pedido, Producto}
\begin{center}
  % Generar con: make diagrams (requiere Java + plantuml.jar)
  % Fuente: images/src_plantuml/04_er_basico.puml
  \fbox{\parbox{0.7\textwidth}{\centering\vspace{1em}[Diagrama PlantUML: ejecutar \texttt{make diagrams}]\vspace{1em}}}
\end{center}
\end{diagramabox}

\begin{lstlisting}[caption={DDL derivado del modelo ER}]
CREATE TABLE clientes (
    id     SERIAL PRIMARY KEY,
    nombre VARCHAR(100) NOT NULL,
    email  VARCHAR(200) UNIQUE NOT NULL
);

CREATE TABLE productos (
    id     SERIAL PRIMARY KEY,
    nombre VARCHAR(100) NOT NULL,
    precio NUMERIC(10,2) NOT NULL CHECK (precio >= 0)
);

CREATE TABLE pedidos (
    id          SERIAL PRIMARY KEY,
    id_cliente  INTEGER NOT NULL REFERENCES clientes(id),
    fecha       TIMESTAMP DEFAULT NOW(),
    estado      VARCHAR(20) DEFAULT 'pendiente'
);

CREATE TABLE pedido_items (
    id_pedido   INTEGER REFERENCES pedidos(id)   ON DELETE CASCADE,
    id_producto INTEGER REFERENCES productos(id) ON DELETE RESTRICT,
    cantidad    INTEGER NOT NULL CHECK (cantidad > 0),
    precio_unit NUMERIC(10,2) NOT NULL,
    PRIMARY KEY (id_pedido, id_producto)
);
\end{lstlisting}

\begin{lstlisting}[caption={Consulta con JOIN múltiple}]
SELECT p.id AS pedido, c.nombre AS cliente,
       pr.nombre AS producto, pi.cantidad, pi.precio_unit
FROM   pedidos p
JOIN   clientes c      ON p.id_cliente = c.id
JOIN   pedido_items pi ON pi.id_pedido = p.id
JOIN   productos pr    ON pi.id_producto = pr.id
WHERE  p.estado = 'pendiente'
ORDER BY p.fecha DESC;
\end{lstlisting}

\begin{entrevistabox}
\begin{enumerate}
  \item ¿Cuál es la diferencia entre notación Chen y Crow's Foot?
  \item ¿Cómo modelas una relación M:N en SQL?
  \item ¿Qué es una entidad débil? Da un ejemplo.
\end{enumerate}
\end{entrevistabox}
